\input{../tex-inputs/latex-leseansicht-vorspann}

\section[Arthur Schnitzler, Jakob Wassermann und Richard Beer-Hofmann an Gustav Schwarzkopf, 5. 8. 1899]{L04433 Arthur Schnitzler, Jakob Wassermann und Richard Beer-Hofmann an Gustav Schwarzkopf, 5. 8. 1899}
\nopagebreak\mylabel{L04433v}
\rehead{ }\normalsize\beginnumbering\briefempfaengerindex{Schwarzkopf, Gustav@\textsc{Schwarzkopf, Gustav}!zzzBeer-Hofmann, Richard@\emph{von Richard Beer-Hofmann}!1899-08-051@{5. 8. 1899}|(be}\briefempfaengerindex{Schwarzkopf, Gustav@\textsc{Schwarzkopf, Gustav}!zzzWassermann, Jakob@\emph{von Jakob Wassermann}!1899-08-051@{5. 8. 1899}|(be}\briefempfaengerindex{Schwarzkopf, Gustav@\textsc{Schwarzkopf, Gustav}!zzzSchnitzler, Arthur@\emph{von Arthur Schnitzler}!1899-08-051@{5. 8. 1899}|(be}
\toendnotes[C]{\smallbreak\pagebreak[2]}
\correspDesc{Versand  durch Arthur Schnitzler, Jakob Wassermann, Richard Beer-Hofmann am 5. 8. 1899 in Altprags
\newline{}Übermittlung  am 5. 8. 1899 in Altprags
\newline{}Zustellung  am 8. 8. 1899 in Wien
\newline{}Erhalt  durch Gustav Schwarzkopf am [8]. 8. 1899 in Wien}\toendnotes[C]{\smallbreak}
\Standort{DLA, A:Schnitzler, HS.1985.1.1897.}
\physDesc{Bildpostkarte, 245 Zeichen
\newline{}Handschrift Arthur Schnitzler: Bleistift, deutsche Kurrent
\newline{}Handschrift Jakob Wassermann: Bleistift
\newline{}Handschrift Richard Beer-Hofmann: Bleistift
\newline{}Versand: 1) Stempel: »\nobreak{}\oindex{Altprags@\textbf{Altprags}|pwk}Altprags, 5. 8. 99\nobreak{}«.   2) Stempel: »\nobreak{}\oindex{I., Innere Stadt@\textbf{I., Innere Stadt}|pwk}Wien 1/1 \textcolor{gray}{1}, 8. 8. \textcolor{gray}{99}, 9–19½V, \textcolor{gray}{Bestellt}\nobreak{}«. }\toendnotes[C]{\smallbreak}\pstart{}{\pb}Herrn \textsc{Gustav Schwarzkopf}\pend{}\pstart{}\textsc{Wien}\oindex{Wien@\textbf{Wien}|pw}\pend{}\pstart{}\textsc{I Tiefer Graben 23}\oindex{Wien@\textbf{Wien}!I., Innere Stadt@\textbf{I., Innere Stadt}!Tiefer Graben 23@\textbf{Tiefer Graben 23}|pw}.\pend{}{\bigskip}
\pstart
           \centering{}{\pb}\textcolor{gray}{\textbf{Der
                  Dürenstein\oindex{Dürrenstein@\textbf{Dürrenstein}|pw} (Pragser Thal\oindex{Pragser Tal@\textbf{Pragser Tal}|pw}).}}\pend
           \vspace{1em}
\pstart
           \noindent{}{\pb}Heute früh um 6 »aufgebrochen«. Um 8, Alt Prags\oindex{Altprags@\textbf{Altprags}|pw}, frühſücken 
      wir und grüßen Sie herzlich. Vorläufig ohne \label{K_L04433-1v}\edtext{Landauer}{\lemma{\textnormal{\emph{Landauer}}}\Cendnote{\textnormal{eine Art Kutsche}}}\label{K_L04433-1}\spacefill\mbox{ArthurSchn}\pend
           
\pstart
           \noindent{}Einige Worte \label{K_L04433-2v}\edtext{Freitag \textsc{Bozen\oindex{Bozen@\textbf{Bozen}|pw}}}{\lemma{\textnormal{\emph{Freitag Bozen}}}\Cendnote{\textnormal{Die Ankunft in Bozen\oindex{Bozen@\textbf{Bozen}|pwk} erfolgte
               erst einen Tag später, am 12. 8. 1899.}}}\label{K_L04433-2}{ }\textsc{\label{K_L04433-3v}\edtext{post. rest}{\lemma{\textnormal{\emph{post. rest}}}\Cendnote{\textnormal{poste restante: postlagernd}}}\label{K_L04433-3}} zu finden, würde 
      \textcolor{gray}{uns} freun.\pend
           \selectlanguage{ngerman}\vspace{1em}\pstart \spacefill\mbox{{[}hs. Wassermann:{]} Wassermann}\pend{}\selectlanguage{ngerman}\vspace{1em}\pstart \spacefill\mbox{{[}hs. Beer-Hofmann:{]} Richard}\pend{}\selectlanguage{ngerman}\endnumbering\briefempfaengerindex{Schwarzkopf, Gustav@\textsc{Schwarzkopf, Gustav}!zzzBeer-Hofmann, Richard@\emph{von Richard Beer-Hofmann}!1899-08-051@{5. 8. 1899}|)be}\briefempfaengerindex{Schwarzkopf, Gustav@\textsc{Schwarzkopf, Gustav}!zzzWassermann, Jakob@\emph{von Jakob Wassermann}!1899-08-051@{5. 8. 1899}|)be}\briefempfaengerindex{Schwarzkopf, Gustav@\textsc{Schwarzkopf, Gustav}!zzzSchnitzler, Arthur@\emph{von Arthur Schnitzler}!1899-08-051@{5. 8. 1899}|)be}\mylabel{L04433h}
\begin{anhang}
\end{anhang}\newcommand{\dateiname}{L04433}\newcommand{\titel}{Arthur Schnitzler, Jakob Wassermann und Richard Beer-Hofmann an Gustav Schwarzkopf, 5. 8. 1899}\newcommand{\editorInnen}{Herausgegeben von Jahnke, SelmaMüller, Martin Anton}\input{../tex-inputs/latex-leseansicht-abspann}
